% Fichier complété par tools/complete_missing_corrections.py.
% Source : DYN/DYN-01/14_Sympact/14_Sympact.tex
% Statut : rédigé (5 question(s)).
\begin{corrigebox}[Corrigé rédigé]
\CorrigeQuestion{1}
Le profil trapézoïdal comporte une accélération constante $+a$ sur
                $t_a$, une vitesse constante sur $t_v$, puis une décélération $-a$ sur
                $t_a$. La vitesse est triangulaire/trapézoïdale et la position est
                quadratique pendant les rampes puis affine pendant le palier.
\CorrigeQuestion{2}
$\boxed{T=2t_a+t_v}$.
\CorrigeQuestion{3}
$\boxed{\omega_{\max}=a\,t_a}$.
\CorrigeQuestion{4}
L'angle total est l'aire sous la courbe de vitesse :
                $\boxed{\Delta\theta=\omega_{\max}(t_v+t_a)=a t_a(t_v+t_a)}$.
\CorrigeQuestion{5}
En combinant la distance imposée et la durée disponible, on résout
                $\Delta\theta=a t_a(T-t_a)$; la racine physique vérifie $0<t_a<T/2$,
                puis $\omega_{\max}=a t_a$.
\end{corrigebox}
